\documentclass[11pt,a4paper]{article}
\usepackage[margin=2.5cm]{geometry}
\usepackage[T1]{fontenc}
\usepackage{lmodern}
\usepackage{microtype}
\usepackage{hyperref}
\usepackage{graphicx}
\usepackage{float}
\usepackage{booktabs}
\usepackage{enumitem}
\usepackage{parskip}
\usepackage{xcolor}
\hypersetup{colorlinks=true,linkcolor=blue,urlcolor=blue,citecolor=blue}
\setlength{\parindent}{0pt}
\setlength{\parskip}{6pt}
\sloppy
\emergencystretch=3em
\hbadness=10000
\begin{document}
\begin{titlepage}
\centering

{\Large University of St Andrews}\\[0.4cm]
{\large School of Computer Science}\\[0.4cm]
{\large Academic Year 2025/26}\\[1cm]

{\large CS5199 Individual Masters Project}\\[2cm]

{\huge\bfseries \textit{A Drop-Aware Engine for an N-Board Chess Variant on a Ring Topology}}\\[0.4cm]

{\Large Infinite Chess Armada}\\[2cm]

\begin{figure}[h]
    \centering
    \includegraphics[width=0.35\linewidth]{sta_logo.png}
\end{figure}

\vfill

{\Large Devansh Chopra}\\[0.3cm]
{\large Student ID: 220026535}\\[1cm]

{\large MSci (Hons) Computer Science}\\[1cm]

{\large Supervisor: Professor Richard Connor}\\[1.5cm]

{\large Submitted to the University of St Andrews}\\[0.2cm]
{\large 18 May 2026}
\end{titlepage}
\newpage
\section*{Abstract}
\addcontentsline{toc}{section}{Abstract}

Infinite Armada Chess, is a multi-board chess variant inspired by Randall Munroe's XKCD comic \#3020 ~[1]. The project introduces a new loop based chess topology where captured pieces move unidirectionally between N connected boards. The human player always controls board 0, while the remaining boards are controlled by Stockfish engines. The main challenge of the system is cross-board decision making - deciding when to drop pieces, where to drop them and which sibling board to request pieces from.

The project is based on John DiIorio’s open-source bughouse repository~[2]. This repository contains a multiplayer frontend for bughouse and a modified chess.js fork that supports reserve pieces and the \texttt{@} notation for drops. The project was developed in two phases. Phase 1 implements a four player two-board bughouse system including server-side rule enforcement,  Stockfish engine opponents and team aware multiplayer logic. Phase 2 extends this into the Infinite Armada a N board variant, currently supporting up to N=10 boards.

The project was tested with a 150 assertion regression test suite as well as engine-strength and pacing tests. In multiplayer testing, the project runs up to N=10 but gameplay is best at N=4 to N=6 levels, where cross-board interaction is meaningful but not overwhelming to the human player. The system was also tested via self-play at N=4, N=6 and N=10 levels.

\newpage
\section*{Declaration}
\addcontentsline{toc}{section}{Declaration}
I declare that the material submitted for assessment is my own work except where credit is explicitly given to others by citation or acknowledgment. This work was performed during the current academic semester except where otherwise stated. The main text of this project report is \textbf{9,758} words long.
In submitting this project report to the University of St Andrews, I give permission
for it to be made available for use in accordance with the regulations of the University Library. I retain the copyright in this work.

\newpage
\section*{Acknowledgements}
\addcontentsline{toc}{section}{Acknowledgements}

Firstly, I would like to thank Professor Richard Connor for his constant support and guidance throughout the semester. I have really enjoyed our weekly meetings throughout the semester where I could to ask for help, discuss my progress and get back on track should I have gone off course. 

I would also like to thank my parents (Dr. Nitin Chopra and Mrs. Radhika Chopra) and brother (Mr. Pranav Chopra) as well as all my family and friends for their support and encouragement throughout the project. 

Lastly, I am thankful to the University and the School of Computer Science for their constant support and understanding. I would also like to acknowledge Randall Munroe, whose XKCD \#3020 comic inspired this project and John DiIorio, whose open source repository laid the groundwork for the project.

\newpage
\section*{Contents}
\addcontentsline{toc}{section}{Contents}
\renewcommand{\contentsname}{}
\tableofcontents
\newpage
\section*{1. Introduction}
\addcontentsline{toc}{section}{1. Introduction}
\subsection*{1.1 The XKCD comic}
\addcontentsline{toc}{subsection}{1.1 The XKCD comic}
Being a chess player I have always known of Randall Munroe's XKCD comic \#3020, Infinite Armada Chess~[1]. When I saw it as a dissertation topic I was immediately drawn in. The comic depicts a chessboard with an endless line of queens behind the normal starting rank for both white and black. The joke lying in the fact that the game can't actually start let alone reach a conclusion. As a chess player, my initial reaction was not to ponder the meaning of "infinite" here, but rather to think about how to make the joke playable. If we take the concept literally and give each side an unlimited number of queens, the game becomes a single tactical puzzle. In this scenario specifically, material advantages and tactics would lose all significance. I came across several online variants that interpreted the comic in this way but none of them approached "infinite" as a topology question, instead they focused on the expanding number of pieces. This project aims to fill that gap.

\begin{center}
\begin{figure}[H]
    \centering
    \includegraphics[width=0.5\linewidth]{xkcd-comic.png}
    \caption{\textit{Randall Munroe’s XKCD \#3020 - Infinite Armada Chess. https://xkcd.com/3020/}}
\end{figure}
\end{center}
\subsection*{1.2 What I built}
\addcontentsline{toc}{subsection}{1.2 What I built}
Rather than giving each side infinite queens, I created my own variant by giving the game an arbitrary number of boards arranged in a circular loop. I split the project into two phases.

\textbf{Phase 1 - Two Board Bughouse System}

Phase 1 began with John DiIorio’s open source bughouse repository~[2], which provided a basic chess UI and some game logic. I extended it into a complete two board bughouse game with real time Socket.IO gameplay, configurable time controls, captured piece transfer, Supabase authentication, ratings, leaderboard and Stockfish engine integration. The general idea is that player 1 (white) and player 2 (black) are 1 team and vice-versa. If player 1 cuts player 3’s black piece, his teammate player 2 can use that black piece in his
turn to place it anywhere on the board with 2 simple constraints which are -
\begin{enumerate}
    \item No direct mate
    \item No direct promotion
\end{enumerate}
This way the players keep exchanging pieces and think of tactics under time pressure
improving their in game discoveries. A variant of this quite recently got termed as Bughouse
Chess.

The main work in this phase was rewriting the rule enforcement layer around Bug.js. All drop rules are checked server side, including no direct mate by drop, no direct promotion by drop and no placement on the opponent’s back rank. This prevents both the client and the engine from bypassing the rules. I also added the official Stockfish engine, so one human can play against three Stockfish opponents at a chosen Elo rating. For Phase 1 I also added a multiplayer option where a room code is generated upon starting a new game which can be shared with friends to join the game room.
\begin{figure}[H]
    \centering
    \includegraphics[width=0.75\linewidth]{phase1.png}
    \caption{Phase 1 - 2 vs 2}
    \label{fig:placeholder}
\end{figure}
\textbf{Phase 2 - Infinite Chess Armada N-Board System}

Phase 2 extends Phase 1 into an N board system, currently configurable from 2 to 10 boards. The human plays on board 0, while all other boards are controlled by Stockfish engines at different Elo levels.

The boards form a circular loop. Pieces captured on board N move into the pocket of board N+1, and the final board feeds back into board 0. Pocket pieces can be legally dropped using the same server side rules from Phase 1. The player can also request pieces from other boards reserves/pocket, although the recommendation logic still needs improvement.

The N=10 limit on boards is a cost constrain, not a theoretical one. Each board needs its own Stockfish subprocess, plus an evaluator so CPU and memory costs grow as more boards are added. Therefore, the comic’s idea of “infinite” is implemented practically as configurable N.

\begin{figure}[H]
    \centering
    \includegraphics[width=0.75\linewidth]{phase2-layout.png}
    \caption{\textit{N=4 boards in Phase 2.}}
\end{figure}

\subsection*{1.3 Why a circular loop?}
\addcontentsline{toc}{subsection}{1.3 Why a circular loop?}
The comic itself is called an armada which means a fleet of pieces. And based on my variant I improvised the metaphor to an armada of boards rather than pieces. I used the concept of a loop since every piece cut on a board 0 is transferred to the pocket of 1, if it was a line board 0 would have no one to feed it pieces and board N-1 would have nowhere to send its captures. Also, because pieces only move one way you can not trade on your board and starve a downstream board or feed it by trading a lot. Nothing in standard chess prepares you for the situation where your opponent's pocket is full because the player two boards back hung a rook.

\subsection*{1.4 Contributions}
\addcontentsline{toc}{subsection}{1.4 Contributions}
Five things, in order of weight.

\begin{enumerate}[leftmargin=*]
\item I added a bughouse rules on top of the \texttt{chess.js}~[10]. \texttt{Bug.js} enforces these rules on the server. 
\begin{enumerate}
    \item A drop which results in mate is not allowed. 
    \item A drop which results in an immediate promotion is not allowed. 
    \item A pawn drop on the first or 8th rank is not allowed. 
    \item A drop on the opponent’s first rank is not allowed.
\end{enumerate}
\item \texttt{InfiniteGame} class is used to manage N boards that are in a loop. For each board it runs a Stockfish process and a shared evaluator. There is a 200ms hardcoded clock between each engine move (since engine vs engine games move extremely fast) and a 250ms engine loop.
\item A drop decision pipeline at \texttt{dropDecision.js} that consists of a search for a mate in one move, an Elo rating drop bias, a heuristic based on candidate drops and a Stockfish comparison of a best drop move against the best move on the board without the drop.
\item A cross board recommendation algorithm. The human can request a piece from any of the teammates boards reserves. The system ranks the boards in terms of how badly the position on each of the boards would deteriorate if it were to give that piece either from the board or the reserves.

\end{enumerate}

The program also consists of FIDE chess rating updates, leaderboard, multiplayer options along with the engines, game notation, time based game (to make it fast paced).

\newpage
\section*{2. Background and related work}
\addcontentsline{toc}{section}{2. Background and related work}
\subsection*{2.1 Chess, bughouse, crazyhouse}
\addcontentsline{toc}{subsection}{2.1 Chess, bughouse, crazyhouse}
For the reader who does not play, regular chess is played between two people on a 8x8 board with each player making alternate moves. The goal is to checkmate the opponent’s king. Bughouse is a 2 vs 2 team variant played on two boards side by side. Pieces captured on one board go into the partner’s pocket on the other board. The partner can then choose to drop that piece as one of their moves. Crazyhouse is the same but played as a 1-vs-1 variant, with captured pieces going into your own pocket instead.

Infinite Armada is somewhere between Bughouse and Crazyhouse~[15, 16, 17, 18]. Infinite Chess Armada is an N board variant played in a loop, with every board having a unidirectional connection to the next. The pocket mechanics are identical to standard bughouse, but now each player has a chain of partners instead of just one.

\subsection*{2.2 The legality of drops}
\addcontentsline{toc}{subsection}{2.2 The legality of drops}
To full proof this variant I implemented the following rules

\begin{itemize}[leftmargin=*]
\item No pawn drop on the 1st or 8th rank. Since a pawn there would imply free promotion or illegal respectively.
\item No drop to directly mate. This is the most important rule. Mate on drop is too would remove most of the skill from the game.
\item No drop for immediate promotion in the same move(Paired with the rank rule)
\item No drop on the opponent's first rank. This would allow unfair advantage to mate quickly.
\end{itemize}
Standard chess rules I follow are from the FIDE Laws of Chess (2023)~[3]. The bughouse conventions come from the ICC~[6] and Lichess~[18] help pages, and my rule set is slightly stricter than Lichess.

\subsection*{2.3 Engines, UCI and why drops are hard}
\addcontentsline{toc}{subsection}{2.3 Engines, UCI, and why drops are hard}
Stockfish is the strongest open source chess engine~[4]. It speaks UCI, the Universal Chess Interface~[11] over stdin and stdout. Therefore, Stockfish does not understand drops normally, which was a big challenge. UCI has no syntax for 'drop a knight on f3'. You can ask Stockfish for the best standard move, but the moment you want the engine to choose between a move and a drop the evalutation heuristic falls drastically.

There is no bughouse engine of any wide use that I can use in place of a normal engine. The older Sjeng engine~[22] is no longer maintained and therefore could not be used. Lc0 (Leela Chess Zero)~[21], the neural net chess engine does not have bughouse weights trained because no one has trained a model of that size for such a special variant. Therefore, I had to use Stockfish for the move side of bughouse and write a drop side engine myself for the times that I need to make a drop instead of a move, which still needs improvement.

\subsection*{2.4 Elo and Glicko}
\addcontentsline{toc}{subsection}{2.4 Elo and Glicko}
Elo is a rating system used in chess and other games~[13]. It updates the ratings of the two players before and after a match by transferring some points from the loser to the winner, where the amount of points transferred is determined by the difference in rating and the expected result of the match (i.e. a highly rated player is expected to beat a lower rated player, thus winning gains fewer points and losing loses more points). The Elo update is typically controlled by a K factor, which for more experienced players can be smaller to reflect that they have played more games.

 For Phase 1 ratings, I used Glicko2 ratings~[7, 14] instead of basic Elo ratings. The Glicko2 ratings are better for new players with very few games, because they take into account the rating deviation in the rating update. For the four player team mode with 2vs2 matches, I couldn’t find any libraries that support rating 2 vs 2 matches in team mode. So I wrote a custom team-Elo function to support this. This code is in \texttt{elo.js} and I tested it with 26 cases in \texttt{tests/elo.test.js}.

\subsection*{2.5 The upstream fork}
\addcontentsline{toc}{subsection}{2.5 The upstream fork}
After approving from my supervisor, I used John DiIorio's open-source \texttt{bughouse} repository~[2] as the starting point. I consisted of a basic frontend, two board Chessground rendering, a Socket.IO setup with \texttt{/lobby} and \texttt{/game} and a chess.js that already handled the reserve data structure and the drop notation. I expanded it with a rule enforcement, Stockfish engine integration, rating system, N board notion, connected supabase and the Phase 2 implementation.

\subsection*{2.6 The rest of the stack}
\addcontentsline{toc}{subsection}{2.6 The rest of the stack}
\begin{itemize}[leftmargin=*]
\item \textit{Chessground }(Lichess) for board rendering and drag and drop~[5].
\item \textit{Socket.IO} for real time channels~[8].
\item \textit{Stockfish 18} as the engine and a shared evaluator~[4].
\item \textit{Supabase} for Postgres and login for multiplayer, replacing the old Postgres plus JWT setup~[9].
\item \textit{Glicko2} npm package for the rating math~[7].
\end{itemize}
\newpage
\section*{3. Requirements and design}
\addcontentsline{toc}{section}{3. Requirements and design}
\subsection*{3.1 Functional requirements}
\addcontentsline{toc}{subsection}{3.1 Functional requirements}

This requirement was split into Phase 1, Phase 2 and shared requirements. All requirements are listed below in priority order of must have (M) or nice to have (N).

\begin{enumerate}[leftmargin=*]
\item  \textit{FR1 (M)} Two board bughouse with four human players.  All moves are validated on the server.
\item \textit{FR2 (M)} Any of the four Phase 1 player slots can be filled with a Stockfish engine with a chosen Elo, so the game runs with 1-4 humans/engines.
\item \textit{FR3 (M)} Captures on one Phase 1 board transfer to the partner’s pocket on the other board in real time. None of the pieces captured on one board are double counted on the other and there are no losses of pieces.
\item \textit{FR4 (M)} A draw on one board does not end the game in draw, a full draw requires all boards to have a tied score.
\item \textit{FR5 (M)} First board to lose, make the entire team lose and ends the game.
\item \textit{FR6 (M)} For all drops in both phases the drop rules are enforced on the server side.
\item \textit{FR7 (M)} Glicko2 ratings calculator and a team Elo update for 2vs2 games.
\item \textit{FR8 (M)} N board (infinite mode) for $N \in \{2, 4, 6, 8, 10\}$ boards with one human on board 0 and rest engines.
\item \textit{FR9 (M)} Transfers of captured pieces in Phase 2 bughouse from one board N to the other board N + 1 real time. All captures made on a board which has already ended are routed to the next active board.
\item \textit{FR10 (M)} Each engine on a Phase 2 board is set to play at a rating of a chosen Elo (800—2500) in matches against other engines. Think time for each engine on a board is set using the pacing formula in Section6.7 for the amount of time there is remaining on the board for that player.
\item \textit{FR11 (M)} The human can ask for a piece from any of the teammate boards pocket in Phase 2. The system will show the human a ranked list of (board, piece) recommendations to request a piece from a teammate and the eval cost of giving up that piece.
\item \textit{FR12 (M)} Team boards in Phase 2 show an eval bar, which updates on every engine move so the human can see what is the status of the rest of his team.
\item \textit{FR13 (M)} The decision whether to drop a piece in a given situation is made by a pipeline consisting of a search for a mate in 1 move by the engine on the human player’s board, followed by a heuristic prune search for candidate moves by the human player on all boards, followed by an Elo scaled bias in favor of a drop for the human player on all boards, followed by a Stockfish evaluated comparison of the best standard move and the best drop move by the human player on all boards.
\item \textit{FR14 (M)} A 60 second reconnect period for reconnecting after a network drop or a laptop lid close in Phase 2 of a game.
\item \textit{FR15 (M)} Time control for a game can be set during game setup. It is set to 5+5 (5 min + 5 sec) as default and a flagging (time out) will end the game on that board with the result.
\item \textit{FR16 (M)} Authentication via Supabase email and password with JWT, ratings and game history in Postgres with row level security on every table.
\item \textit{FR17 (M)} All move legality checks happen on the server, the client cannot submit an illegal move and the engine cannot bypass the rule layer because both code paths share the same predicates.
\item \textit{FR18 (N)} Multiplayer functionality via room codes so friends can join from the same URL without an explicit invite.
\item \textit{FR19 (M)} Seat management in the lobby waiting room so that players can swap seats and add engines until all four are filled.
\item \textit{FR20 (N)} Profile auto creation on signup via the \texttt{handle\_new\_user} trigger.
\item \textit{FR21 (N)} Local game mode that bypasses Socket.IO entirely for offline play.
\item \textit{FR22 (M)} A post game update that shows the result, rating deltas (Phase 1) and a final state summary (Phase 2) before the players return to the lobby at the end of a game.
\item \textit{FR23 (N)} A leaderboard page that lists the top Phase 1 players by current rating.
\item \textit{FR24 (M)} Leaderboard is updated after each game based on the calculations.
\end{enumerate}

\subsection*{3.2 Non-functional requirements}
\addcontentsline{toc}{subsection}{3.2 Non-functional requirements}

\begin{enumerate}[leftmargin=*]
\item \textit{NFR1} Engine move latency under 200ms.
\item \textit{NFR2} In Phase 2, the engine boards will never move faster than a human at the same time control. This is enforced by \texttt{ENGINE\_MIN\_GAP\_MS = 1500} and is removed when engine has less than 1min.
\item \textit{NFR3} The drop decision pipeline must return a move within the think time set (300ms to 1500ms) depending on remaining clock.
\item \textit{NFR4} The server validates, broadcasts and only then does the client render moves.
\item \textit{NFR5} Memory scales linearly with active games and with $N$. An $N{=}10$ game uses 11 Stockfish processes (10 boards plus 1 shared evaluator).
\item  \textit{NFR6} Row-level security on every Supabase table; JWT validation on every Socket.IO event; service-role access restricted to the server's admin client.
\item \textit{NFR7} Rule tests are deterministic but engine strength tests are stochastic and must clear thresholds set (stronger engine wins >= 3/5, weaker wins >= 1/5).
\item \textit{NFR8} Code organisation of the rules (\texttt{bug.js}), the engine (\texttt{stockfishEngine.js}), the drop decision maker (\texttt{dropDecision.js}) and the evaluation (\texttt{bughouseEval.js}) are organised in completely independent layers of the system. Changes to any of the above should never affect the code in the other layers.
\item \textit{NFR9} Phase 1 and Phase 2 share the rule layer but run on separate Socket.IO namespaces, so a bug in one cannot affect the other phase.

\end{enumerate}
\subsection*{3.3 Constraints}
\addcontentsline{toc}{subsection}{3.3 Constraints}
Due to using the bughouse repository, I had to use React 15~[12], Webpack 3 and Babel 6. The software uses Node 20, as that is what the Supabase client expects. Furthermore, Stockfish has no built in multi instance API so each board gets its own process. The drops are outside of UCI, so had to add custom code which still needs to be improved. Additionally, there is no bughouse training data of sufficient scale to train a neural net chess engine, even if I had the time to do so.

\subsection*{3.4 Key design choices}
\addcontentsline{toc}{subsection}{3.4 Key design choices}
Four key decisions that I made during this project
\begin{enumerate}[leftmargin=*]

\item \textit{Engine integration -}  I used Stockfish and add a drop layer above it. I could not build my own engine due to the lack of games for a bughouse variant. I used Stockfish as it is the best chess engine and adding my drop logic on top of it correctly makes it account for pocket pieces as well.
\item \textit{Drop decision -} I was deciding between a pure heuristic, ML model or heuristic plus Stockfish eval. I implemented the third by the end as no bughouse training data, so ML was not possible. A pure heuristic would perform poorer in comparison to a Stockfish for a standard move. Therefore, the hybrid uses Stockfish for the move side and my eval module for the drop side.
\item \textit{Multi-board topology -} Before starting the project I discussed these 3 possibilities with my supervisor (line, loop or graph).  I decided to use a loop since this way I avoid dealing with special end of line positions where the last board can not transfer pieces to the first board.
\item \textit{Engine pacing -}  Another issue I was facing was that the engines would blitz during a game and I could not keep up with their speed. This would lead to me have 10 pieces in my reserve even though I have only played 5 moves on my board. Therefore, I scaled the clocks, that is at low clocks they would need to actually try to flag. The formula to scale the think time by the clock is \texttt{clamp(remainingMs/25, 300, 1500)} for less than 30 seconds of remaining time. For more than 30 seconds of remaining time, a fixed think time of 1500ms is used.
\end{enumerate}
\subsection*{3.5 What changed from the DOER plan}
\addcontentsline{toc}{subsection}{3.5 What changed from the DOER plan}
There was one major change from the DOER plan, instead of writing out my own evaluation of how chess pieces move and positions should be evaluated, I used stockfish for the standard chess moves and add a layer on top to account for piece drops. Originally the engine was a big part of the dissertation but after speaking with Professor Connor I realized that making an engine was not realistic for the scope of this project given the time and available dataset constraints for a bughouse variant.

\newpage
\section*{4. Software engineering process}
\addcontentsline{toc}{section}{4. Software engineering process}
This project was developed over a span of 11 weeks, with each having a weekly checkpoint meetings with the supervisor. Where we discussed the progress, what should be done next and will it fit well into the variant. Every week a new feature (small enough to be working within a week or two) along with the necessary tests for testing the new feature was delivered. The version control for this project was handled via GitHub.

I split this project into two phases. Phase 1 focused on building and validating the rule system by testing and manual playtests before moving onto Phase 2, which extended the system into the infinite N board topology. Each phase was manually tested before moving forward. The original DOERs proposal was modified over the course of this work. The largest change from the original plan was that instead of writing a full evaluation engine from scratch, the Stockfish chess engine was used and a drop decision pipeline was added on top.

Additionally for testing I have used test-driven development (TDD) for the rule system in \texttt{FordropRules.js} as well as the team Elo system in \texttt{elo.js}, while the rest of the system I have tested in an iteratively. There is a test suite of 150 tests in 5 files, which can be run with \texttt{npm test}.

VS Code was used for editing, ESLint (set to the same configuration as the upstream project) was used for style checking and a single \texttt{run.sh} script was written to set up a development environment (installs  Node 20, Stockfish and Python dependencies for the other part of the project,and then starts up the local server).
\newpage
\section*{5. System architecture}
\addcontentsline{toc}{section}{5. System architecture}
\subsection*{5.1 Three layers}
\addcontentsline{toc}{subsection}{5.1 Three layers}
The architecture is a 3 layer stack. On the client side a browser loads a single page React 15 app as a webpack bundle by the server. The Node server (Express plus Socket.IO) holds the game data in memory and starts one Stockfish subprocess per board of a game, plus one shared evaluator for all the boards of that game. Supabase is used to manage Postgres plus auth. Stockfish uses UCI over stdin and stdout.
\begin{figure}[H]
    \centering
    \includegraphics[width=0.5\linewidth]{system_topology.png}
    \caption{Enter Caption}
\end{figure}

\subsection*{5.2 Frontend}
\addcontentsline{toc}{subsection}{5.2 Frontend}
The application uses React 15 with Redux 3 and react-router. The client has 5 major files -
\begin{enumerate}
    \item \texttt{AuthPage.jsx} which is for logging in and out of the server
    \item \texttt{LobbyPage.jsx} joining games using room code
    \item \texttt{LobbyWaitingRoom.jsx} for configuring the lobby and multiplayer games
    \item \texttt{InfiniteSetupPage.jsx} for setting up games upto N=10
    \item \texttt{InfiniteGamePage.jsx} for the Phase 2 Infinite Armada gameplay.
\end{enumerate}

For Phase 1, I reused DiIorio’s \texttt{GameBoardsComponent.jsx} with the added sidebars for actions, moves and game information. For Phase 2, I created a custom centre-board plus mini-boards layout for the client. The main board is rendered with Chessground 6.5.8 and the sibling boards are rendered with a very lightweight SVG component in MiniBoardSvg.jsx. I did not use multiple Chessground instances for the mini-boards because running up to nine extra boards was too slow on my laptop. The smaller boards did not need to have drag and drop interaction so I created a very lightweight SVG component instead.

The Redux store for the client is structured by feature (game, user etc) with flat reducers. The changes for in game state occur on the server therefore the client does never apply a move before it has been validated by the server. I also created a separate \texttt{Clock.js} for the rendering of the clocks on the client. In the original fork the clock logic was mixed into the Redux reducer and behave incorrectly upon reconnect.

\subsection*{5.3 Backend}
\addcontentsline{toc}{subsection}{5.3 Backend}
The backend consists of an Express app located in \texttt{src/server/app.js} which exports \texttt{/api/lobby}, \texttt{/api/ratings}, \texttt{/api/infinite}, \texttt{/api/games}, \texttt{/api/auth}, \texttt{/api/users} and \texttt{/api/leaderboard}. There is also a SPA catch all route for non-API GET requests which returns \texttt{index.html} so deep links work. Also, all static files (the Webpack bundle etc) are served from \texttt{src/client/} by the static middleware.

The three Socket.IO namespaces are \texttt{/lobby}, \texttt{/game} and \texttt{/infinite}. These were inherited from the original project and Phase 2 added the \texttt{/infinite} namespace. Because of the events in Infinite Armada like \texttt{request\_piece}, \texttt{state}, \texttt{eval\_update} etc, these do not need to coexist with the events from Phase 1, like \texttt{offer\_draw}, \texttt{accept\_resign} etc, and thus keeping them separated made the logic much simpler. Bugs in one namespace would not affect the other.

I chose to store active \texttt{InfiniteGame} objects in memory using a \texttt{games} map indexed by game ID. Since, each infinite game contains multiple Stockfish subprocesses, timers and active UCI engine state, so persisting all of them would require significantly more infrastructure than this project needed. Therefore, active games are lost if the server restarts but since most games at 5+5 time controls last under 25 minutes, this was an acceptable trade off for the project.

\subsection*{5.4 Database and auth}
\addcontentsline{toc}{subsection}{5.4 Database and auth}
I used Supabase for both authentication and the Postgres database, with row-level security enabled on all tables. I made this choice, since I plan to deploy this on Vercel for multiplayer functionality. The schema contains \texttt{profiles}, \texttt{games}, \texttt{rating\_history} and the Supabase managed \texttt{auth.users} table. A Postgres trigger automatically creates a profile row whenever a new user signs up, avoiding race condition issues from earlier versions that required separate profile creation requests.

The \texttt{games} table uses JSONB fields, \texttt{engine\_slots} and \texttt{engine\_levels}, to support human and engine player setups. The \texttt{rating\_history} table stores per game rating updates with indexed queries for leaderboard access. Authentication uses Supabase JWT tokens. The token is stored in \texttt{localStorage} by the client and is included in Socket.IO events, and the server verifies it before processing requests.

\begin{center}
\fbox{\begin{minipage}{0.85\textwidth}\centering
\textit{[Figure 4.4 to be inserted]}
\par\vspace{0.4em}
Version-stamped event flow. Client A makes a move; the server bumps version 41 to 42; both clients apply. Client B drops; the server bumps 42 to 43 to 44 in the meantime. Client B reconnects; the server sends full state v44 and B is back in sync.
\end{minipage}}
\end{center}
\newpage
\section*{6. Implementation}
\addcontentsline{toc}{section}{6. Implementation}
\subsection*{6.1 Phase 1 -  2 vs 2}
\addcontentsline{toc}{subsection}{6.1 Phase 1 -  2 vs 2}
For Phase 1 of this project I developed a 2 vs 2 bughouse variant with two players per team playing on two boards side by side, where pieces that are captured on one board are dropped in the pocket of the partner on the other board. The game and its rules are well known in the online chess community, therefore this implementation only consists of a way to enforce the rules and to integrate an engine into the games for the human players missing.

 The state model is implemented in \texttt{Game.js}. It extends DiIorio’s repo of the state model for bughouse chess to support multiple engine slots and to properly keep track of the rating of all four players. A Game object contains 4 player IDs (one for each slot), 4 clocks (only ticking on the turn of the player whose board the clock is for), 2 FENs (Forsyth-Edwards Notation) chess notation for the 2 boards, 2 reserve arrays per color per board, 1 list of promoted pieces per board, and team-level draw and resignation status. The team relationship is decided by the position of the slots i.e slots 0 and 3 form a team (white on board A and black on board B), and slots 1 and 2 form the other team. Thus the partner of slot 0 is slot 3, and captures by slot 0 go to the pocket of slot 3 on the other board.

The slot of a player is assigned in the lobby. The \texttt{LobbyPage.jsx} shows all open slots (for both users and engines) when a new game is created along with a room code. The \texttt{LobbyWaitingRoom.jsx} tile is assigned to a human player (profile name). The other three are assigned to is the waiting room for the other players or engine before the game starts. The Elo of an engine can be chosen between Beginner and Grand Master. A room owner can switch the slot of another player, kick an engine from a slot or take the slot of an a human player and so on. A game is ready to start when all four slots are filled. The slot configuration is stored in two JSONB columns, first  \texttt{engine\_slots} which consists of the slots which are engines and second \texttt{engine\_levels} which is the Elo for each engine. The slot configuration of a game is always read together with the rest of the game and is never queried in isolation. Therefore I chose JSONB instead of a separate table \texttt{game\_slots}

\begin{figure}[H]
    \centering
    \includegraphics[width=0.75\linewidth]{phase-1-lobby.png}
    \caption{Game Lobby}
\end{figure}

The Phase 1 implementation for engines is in \texttt{engineManager.js}. It differs from the  Phase 2 in terms of the unidirectional loop and partner position awareness. The function \texttt{computePartnerNeed} computes how much a certain captured piece would help the partner. For example the engine will be biased to capture a knight if your partners opponent has just castled and a fork is possible using the knight. The two constants \texttt{LAMBDA = 0.3} and \texttt{MU = 0.25} are used to weight the partner-needs and threat from opponent. These two constants were hand tuned after playing multiple games at the interim demo stage, so they are not perfect and can be improved using a ML model, but they make the Phase 1 engine play like a good bughouse partner.

\subsection*{6.2 Rules Applied}
\addcontentsline{toc}{subsection}{6.2 Rules Applied}
The rule layer is in \texttt{bug.js} and \texttt{updateGame.js} which accepts the move on behalf of the server for all game play, and checks four different rules for each drop, all server-side. 

These four rule checks are \texttt{isPawnRankDrop()} which prevents pawns dropping to the player’s 1st or 8th rank, \texttt{isBackRankDrop()} which prevents any piece dropping to the opponent’s 1st rank, \texttt{isPromotionOnDrop()} which prevents any pawn drops that immediately result in promotion and \texttt{isCheckmateOnDrop()} which creates an independent FEN clone, makes the proposed drop and returns false if the resulting board position is a checkmate.

The checks for illegal moves in \texttt{updateGame.js} are executed prior to actually updating the game state. If any of the checks fail then the server will emit an \texttt{illegal\_move} event to the user and the game state on the server remains unchanged. This was important because the path that the Stockfish engine takes through the system does not involve the react client and therefore both human and engine moves would need to use same set of rules to be validated on the client.

The most tricky check to make was the check for checkmate on drop. In normal human chess games this check is only performed a few times per game of chess. However, in bughouse with 8 pawns per player for each color, hundreds of drop possibilities may be considered in a short amount of time in Phase 2. To try to keep the time required to make this check as low as possible, the FEN clone used by this check is reused between calls. This is one of the few places in the system where such an optimization works because N in this case can be max 10 boards and there are multiple engines that could be making the same check call in parallel. And in the worst case scenario in a game of 10 boards an engine/human could have 10x8 =80 pawns to drop.

\subsection*{6.3 Phase 1 socket flow and team awareness}
\addcontentsline{toc}{subsection}{6.3 Phase 1 socket flow and team awareness}
Phase 1 of the server client uses the Socket.IO \texttt{/game}. On the client we send move, offer draw, accept draw, decline draw and resign. On the server we send board\_update, clocks, offer\_draw\_received and game\_over messages. The important bit is how team aware the server is for these messages.

A draw on one board means the entire game is draw. The game ends and a score of 1/2 - 1/2 is assigned to both team.
Resignation on the other hand is also done team wise. If a player either loses or resigns then the score is set to 0-1 for the entire game and all games end.

The \texttt{board\_update} in a single message sends both boards FENs, reserve arrays and one version bump. This helps to keeps captures and reserve updates atomic. If the updates were split into two events a slow client might see one board updated while its partner’s pocket is still out of sync. Therefore,using one event gives one consistent render.

\subsection*{6.4 Phase 2 - N-board loop}
\addcontentsline{toc}{subsection}{6.4 Phase 2 - N-board loop}

In \texttt{infiniteManager.js}, every \texttt{InfiniteGame} has $N$ boards in a circular list. Each board stores its FEN, white and black reserves and clocks, latest engine evaluation, terminated flag. There are two main intervals that manage the game - a clock interval every 200ms that ticks the active player's clock on every board in progress, and an engine loop that runs \texttt{runEngineLoop} once per cycle for every engine board.

The boards are arranged in a circular list and therefore all boards are connected in a loop. The connection is done in \texttt{syncReservesAndFlow}. In every move, the pieces that been captured are put in the reserve of the next board in the list. Since the boards are in a circle, we connect to the next board by \texttt{(board.idx + 1) \% numBoards}. The \texttt{\%} turns the connection into a closed list, and the rest of the boards are connected by normal indexing.

Boards that have finished (by checkmate or flag) are considered dead boards, but are still passed captured pieces from the boards above it (to prevent their pieces from vanishing into space). For phase 2, a sibling board finishing does not end the game as a whole (i.e. the game ends when the human player’s board finishes, i.e. board 0).

\begin{figure}[H]
    \centering
    \includegraphics[width=0.75\linewidth]{n-board-ring.png}
    \caption{N-board ring flow}
\end{figure}

\begin{center}
\fbox{\begin{minipage}{0.85\textwidth}\centering
\textit{[Figure 5.1 to be inserted]}
\par\vspace{0.4em}
The N-board ring at N=4. Four board boxes; one-way arrows from 0 to 1 to 2 to 3 to 0. Human icon on board 0, engine icons on 1, 2 and 3.
\end{minipage}}
\end{center}
\subsection*{6.5 Engine integration}
\addcontentsline{toc}{subsection}{6.5 Engine integration}
\texttt{\texttt{StockfishEngine}} in \texttt{stockfishEngine.js} is a simple wrapper around \texttt{spawn('stockfish')} for the normal UCI~[4, 11] messages \texttt{uci}, \texttt{uciok}, \texttt{setoption}, \texttt{isready}, \texttt{readyok}. Engine strength is determined by \texttt{UCI\_LimitStrength} and \texttt{UCI\_Elo} for Elo greater than 1320. For lower Elo ratings, \texttt{UCI\_LimitStrength} is disabled and Elo is mapped roughly to Skill Level, using \texttt{Math.floor((elo * 20) / 1319)}.

UCI is stateful that is stdin/stdout stream. So even with a single call to getBestMove or evaluation, it is it possible that the output will get corrupted by a prior call to getBestMove or evaluation on the same engine. To avoid this problem, all of the calls on an engine are serialized through a single chain, which is stored in the \texttt{busy} field of the engine. This way, the requests for the engine are processed one at a time. This solved a previously timing bug where a clock tick would occur and cause UCI text from the two requests to get mixed up. Also, the failure handling for UCI for engine requests which either timed out or failed now cause the engine to be marked as not started. The engine then checks for \texttt{started} before each move if it is false, tries a single respawn otherwise falls back to a random legal move (using Bug.js). This is to prevent the game from getting stuck in a deadlock, it is not intended to be good moves, rather to keep the board alive.

\subsection*{6.6 The drop-decision pipeline}
\addcontentsline{toc}{subsection}{6.6 The drop-decision pipeline}
I came up with the drop pipeline structure based on my personal experience of finding the best attack on the board. It works in five steps

\begin{enumerate}
\item \texttt{Check for mate -} Checks for a mate in one by dropping any piece in the pocket for checkmate.

\item \texttt{Strong drop options -} If there is no checkmate drop available, the program searches for a eval of the strongest legal drops in the given position (e.g mate in two). The number of such moves retrieved by the program is kept short in order to be able to process them quickly in the following steps. This short list of the best legal drops is generated by the program by using simple heuristics of chess, e.g king pressure, checks, forks etc.

\item \texttt{Elo behaviour -} The engine uses the chosen Elo to determine how impulsive the engine should behave. For lower rated engines the engine will play the best looking drop immediately, whereas a stronger engine will evaluate all options.

\item \texttt{Compare with Stockfish -} It calculates the best normal move with Stockfish and then compares it to the best possible drop with the short list created in step 2. The move which puts the opponent in the worst position is chosen by the system.

\item \texttt{Use fallback if needed -} If Stockfish fails to return a move or no strong drops were found, we play a random legal move to prevent the board from freezing.

\end{enumerate}

\begin{center}
\fbox{\begin{minipage}{0.85\textwidth}\centering
\textit{[Figure 5.3 to be inserted]}
\par\vspace{0.4em}
Drop-decision pipeline. Five boxes: legal candidates -> mate-in-one search -> heuristic prune -> Stockfish move plus drop comparison -> pick-or-fallback. Side branches for the Elo-bias and random-fallback exits.
\end{minipage}}
\end{center}

\subsection*{6.7 Real time sync}
\addcontentsline{toc}{subsection}{6.7 Real time sync}
The server sends down four events for the client can handle which are \texttt{state} (a full context of all the boards), \texttt{board\_update} (a single board has changed), \texttt{clocks} (just the clock fields have changed) and \texttt{eval\_update} (the eval percent for a single board has changed). All of the events have a single field, all carry the monotonic version. The reconnect path is critical here. When a client re-connects to the server, it re-sends a \texttt{join\_room} event for the stored gameId and the client’s JWT. The server clears any pending disconnect timer for that user and sends a full state. The client replaces it so a tab that has missed an unknown number of events ends up correct.

\subsection*{6.8 The drop UI/UX}
\addcontentsline{toc}{subsection}{6.8 The drop UX}
The drop UI has 3 states \texttt{idle}, \texttt{drop-armed} and \texttt{drop-placing}. In idle mode, clicking a reserve piece activates it for a drop on the board legally. In drop-armed mode, clicking a square will drop the piece there if legal. drop-placing mode is activated briefly when a piece si being placed then returns to idle again.

Phase 2 adds a new feature \texttt{request piece} on the right side. It shows all the boards and their reserves along with the best board to request a piece from based on the eval bar. Selecting a piece from there would remove it form the reserves on that board and add it to yours. I thought this feature was interesting to break the unidirectional loop only for the human who is playing against highly rated stockfish engine. If I had more time I would have liked to work on the logic for this further.

\begin{center}
\fbox{\begin{minipage}{0.85\textwidth}\centering
\textit{[Figure 5.5 to be inserted]}
\par\vspace{0.4em}
The sibling-board recommendation panel. Ranked list of (board, piece) pairs annotated with the eval delta if the piece is removed from the source board.
\end{minipage}}
\end{center}
\subsection*{6.9 Layout}
\addcontentsline{toc}{subsection}{6.9 Layout}

For Phase 1, a normal bughouse simple 2 board view is shown in the center along with the FEN notation on the left. The reserves are next to the respective users clock. The user can analyse by using the back and forward arrows to go back in the game (non editable) and see the position.

\begin{center}
\fbox{\begin{minipage}{0.85\textwidth}\centering
\textit{[Figure 5. to be inserted]}
\par\vspace{0.4em}
Image of back and forward arrows
\end{minipage}}
\end{center}

For Phase 2,  the center has the human board with all the other boards shown at the top along with the eval bar. The user can click on any board to select it as the main board to view. The request piece feature is on the right which shows the reserves per board.

Earlier I considered selecting a number on the keyboard to open that board, but I rejected the idea since cross board awareness is the strategic point of this variant. Therefore, hiding the boards would have not have the same effect on the users mentality while playing and even asking for pieces.

\section*{Ethics}
\addcontentsline{toc}{section}{Ethics}

This project was reviewed by the School of Computer Science Ethics Committee and approved under reference \textbf{CS15727} as a generic artefact evaluation. The risk category was assessed as minimal.

\subsection*{Participants}
\addcontentsline{toc}{subsection}{Participants}
The only human involvement other than me was a small set of voluntary playtest sessions with University of St Andrews students (mostly fellow chess players), six 4 player bughouse sessions were played for Phase 1 and self play feedback for Phase 2.

\subsection*{Data Collected and Retention}
\addcontentsline{toc}{subsection}{Data Collected and Retention}
The system only collects the information needed to play the game. This includes an email address for Supabase login (users can also enter as guest if username is available), a username chosen by the player and game related data such as moves, clocks, FEN positions and rating changes. Feedback from playtesting was stored as anonymous text notes. All saved data is stored in a Supabase Postgres database protected using row-level security and JWT authentication for every Socket.IO event. Phase 2 infinite mode games are stored only temporarily in server memory and are deleted once the game ends. Most of the players requested to have their account removed from the database after play testing.

\subsection*{Compliance}
\addcontentsline{toc}{subsection}{Compliance}
The project follows the University of St Andrews Data Protection and Good Research Conduct guidelines. The ethical approval letter (reference CS15727) is on file with the School.

\newpage
\section*{7. Testing and verification}
\addcontentsline{toc}{section}{7. Testing and verification}

The tests can be run in the project root using \textbf{npm test}. Full test summary in \texttt{Appendix B}

The testing setup (\texttt{tests/test-utils.js}) contains helper functions for testing. Each test file tracks its pass and fail count and gives the result in the end. The \texttt{tests/run-all.js} runs each file as a separate node process, reads the reuslts and if any fail it exits with and error. 

\subsection*{7.1 Coverage}
\addcontentsline{toc}{subsection}{7.1 Coverage}

I have added 150 tests in 5 test files.  \texttt{bughouse.test.js} and \texttt{elo.test.js} test the Phase 1 rule and team rating logic. \texttt{dropDecision.test.js},\texttt{engineStrength.test.js} and \texttt{infiniteMode.test.js} test Phase 2 implementation including engine pipeline and N board system. Tests for both the phases were written simultaneously during respective phase.

\begin{itemize}

\item \textit{\texttt{bughouse.test.js}} - 82 tests to check legality of drops, occupied square checks, king safety, reserve duplication, capture transfer, undo and FEN restoration, check by drop, tactical drop tests, move generation tests, evaluation check, edge cases such as en passant, castling and not enough material with reserve pieces (for 50 move draw).
\item \textit{\texttt{dropDecision.test.js}} - 21 tests that check for mate in 1 drops, correct ordering of user drops, applying of drops to FEN positions, illegal drops and correct move selection by the engine.
\item \textit{\texttt{elo.test.js}} - 26 unit tests for \textit{\texttt{elo}} function updating team Elo after wins, losses, draws as well as edge cases like K factor scaling within bounds, rating floors, missing profiles for updates and returning the updated rating.
\item \textit{\texttt{engineStrength.test.js}} - 2 assertions, skipped if Stockfish is unavailable. It verifies that the stronger engine performs well overall. For this, I added 5 curated positions comparing a Elo 1300 to a Elo 2500 engine.
\item \textit{\texttt{infiniteMode.test.js}} - 19 tests, skipped if no Stockfish installed. These test game initialization, the two engines playing until one of them is checkmated, pacing, reconnect, human moves, engine drops in forced mate position and game termination.

\end{itemize}

\begin{center}
\fbox{\begin{minipage}{0.85\textwidth}\centering
\textit{[Figure 6.1 to be inserted]}
\par\vspace{0.4em}
\texttt{npm test} output, all five files passing with assertion counts and the runner's summary line.
\end{minipage}}
\end{center}
\subsection*{7.2 What is not tested}
\addcontentsline{toc}{subsection}{7.2 What is not tested}
First, I manually tested Chessboard rendering and React connections therefore did not add any frontend tests. Secondly, I was not able to add any load tests. Although, I manually tested 2 simultaneous N=10 games which created 22 Stockfish processes which worked fine. Thirdly, although I dont have tests for socket reconnects. The reconnect system worked well in manual testing, but I did not intentionally break TCP connections during games to test failure behaviour. Given the scope of the project and the time constraint, I thought testing the engine, chess rules and the overall infinte armada superseded these. However, I have manually tested the system and have fixed console errors as well.

\subsection*{7.3 Bugs found during testing}
\addcontentsline{toc}{subsection}{7.3 Bugs found during testing}

The following three bugs where identified during the testing of Bug.js. The first bug was in \texttt{infiniteMode.test.js}. In some cases the engine would checkmate the human player but the game would fail to end properly. This issue was caused by the fact that \texttt{checkBoardEnd()} is a function that is called after a successful move. In some of the mate in one cases for the engine’s fallback path the check for a mate in one move would be skipped. Thus the function for checking the end of the board for check would not be called. Fixed it by placing the function \texttt{checkBoardEnd()} at the start of the function \texttt{makeEngineMove()}.

Another issue was found during manual testing of  remove pawn on promotion (pawn gets removed after promoting to queen). Initially setReserves only copied the reserve arrays and not the piece objects within them. This meant after undoing after such a move, the promoted piece would still be there on the board. This was fixed in Phase 2.

The third bug was the Engine vs Engine speed of moves. They were progressing too fast and my reserves where getting filled with pieces even though I would have not traded anything. This was fixed by correctly mapping \texttt{ENGINE\_MIN\_GAP\_MS} with \texttt{lastMoveAt} instead of \texttt{lastTimestamp} so that the delay in moves keeps the speed on both the boards even.

\newpage
\section*{8. Evaluation}
\addcontentsline{toc}{section}{8. Evaluation}
\subsection*{8.1 Engine strength}
\addcontentsline{toc}{subsection}{8.1 Engine strength sanity}
As mentioned above in 6.1, I added an engine strength test to set up two Stockfish engines of Elo 1300 and Elo 2500 and for each of them run five test positions (200ms between moves). The test positions are an Italian Opening position, a king and pawn endgame, a position with a developed knight to attack, a position where a centre pawn can be captured and an endgame position with advantage. The 2500 Elo engine finds the best move for 4 out of 5 test positions that had been curated for this test. The 1300 Elo engine finds the best move for all 5 test positions, which is not surprising as the test positions chosen were all very simple and had been chosen for that reason. The purpose of this test is not to do Elo calibration itself (that would require a hundred or so games each for Bughouse and Chess) but to check that the engine starts up, talks UCI correctly and returns sensible moves in simple positions.

\subsection*{8.2 Pacing}
\addcontentsline{toc}{subsection}{8.2 Pacing}
I tested the engine speed by not moving on an N=2 game and observing the idle board showing that the engine moves only. I checked the move timings every 100ms over the course of 7 seconds. With a pacing limit of 1500ms the moves are spaced about 1.2 seconds apart which seems to be fast enough to keep the games readable but not so fast that the boards move too quickly.

For an application like chess self play testing was a much more productive approach. For an N=4 game with a 5+5 time control, it was completely possible to follow and understand all the other boards and it felt natural to jump to another board to assess the position and request a piece if necessary. However, at N=10, even though individually each board moved at a reasonable pace, it was not possible to follow all of them in sufficient detail.

\begin{center}
\fbox{\begin{minipage}{0.85\textwidth}\centering
\textit{[Figure 7.2 to be inserted]}
\par\vspace{0.4em}
Engine move-gap distribution on a sibling board across a 10-minute observation. Buckets: under 1.2s, 1.2 to 2s, 2 to 3s, over 3s. Most of the mass falls in the 1.2 to 2s bucket.
\end{minipage}}
\end{center}
\subsection*{8.3 Playability at N = 4, 6, 10}
\addcontentsline{toc}{subsection}{8.3 Playability at N = 4, 6, 10}
I played 10 games on each N, all games 5+5, all engines were set to Elo 1500. For each game, I scored them on a combination of 3 criterias - 
\begin{enumerate}
    \item Did the request piece across boards matter  (i.e Did pulling a piece from a sibling board change the outcome on board 0? 
    \item Did the sibling boards remain tactically meaningful or did they drop into complete noise? 
    \item Did game end in a nice clean fashion (i.e. no deadlocks)?
\end{enumerate}

\begin{itemize}[leftmargin=*]
\item \textit{N=4.} The cross board request mattered in 7 out of 10 games. The sibling boards stayed tactically meaningful in 9 out of 10 games. All games ended cleanly with no deadlocks.
\item \textit{N=6.} The cross board request was critical in 5 out of 10 of these games. The sibling boards were meaningfully close in 8 out of 10 games. All games ended cleanly with no deadlocks.
\item \textit{N=10.} The cross board request mattered in 3 of 10 games. The sibling boards were tactically meaningful in 4 of 10 games. All ten ended cleanly.
\end{itemize}

The variant is best played at N=4 to N=6. N=10 is also perfectly fine from a technical standpoint (no crashes, no deadlocks, correct pacing). But the strategic cross board loop gets severely weakened. In most of the games at N=10, it takes too long for pieces to propagate from board 4 where they were moved to my pocket on that board. In most of the games I won or lost on board 0 without ever having pulled a piece from a sibling board. This result taught me that the interesting strategic playing happens in the middle, not at the edges. I would like to work on the cross board loop so that engines can also ask for a piece and save a weak position.

\subsection*{8.4 Compared to bughouse and to parallel chess}
\addcontentsline{toc}{subsection}{8.4 Compared to bughouse and to parallel chess}

Compared to bughouse, my multiplayer (basic mode) and N=2 (Infinte Armada) version are similar, but the piece transfer strategy is different. In bughouse, my partner can hold onto a piece and wait until I require it, this stops the flow of pieces and can sometimes become overwhelming if my opponent gets 6 pocket pieces in one go. In my variant,  I fixed this by automatically transferring pieces to the partner both in the multiplayer and Infinite Armada. 

Compared to parallel chess, where multiple normal chess games run in parallel without any piece transfer. To test this, I added a patch to disable flow state and tested it with two N=4 games. The experience was like tournament style chess where each game is independent without any piece flow but counts towards the overall team score. In this format you usually focus on your own board and rarely look at the other boards just to get an update. This is not an engaging enough experience for my variant. But it gave me a great idea, in the future I would like to expand the software to have a tournament mode as well where people and engines can compete over a week. This would although require a pairing software as well to update the ratings and match each player for the next round based on official FIDE rules.

One of the N=4 game that I played was quite interesting, I was down a pawn on board 0 and in a weak position. My teammate (engine) on board 2 traded rooks, which landed in board 3 pocket. The system recommended  me that it was the best piece to request, since board 3 was losing (eval bar dropped) and board 1 was winning. I got the rook, placed it on my board and exchanged material to gain advantage and won. Overall, this game was a great testing experience and see all the board interact with each other.

\subsection*{8.5 Phase 1 playability and analysis}
\addcontentsline{toc}{subsection}{8.5 Phase 1 playability and analysis}

To test this, I ran 6 sessions of 4 player (3 sessions at 5+5 and 3 sessions at 3+2). Two of the players were serious players and the other four were casual players.

\begin{itemize}

\item The rule set worked correctly in all 24 games. All difficult mate by drop situations were handled correctly as well.

\item The users quite liked the draw offer functionality. In one of the sessions a tester even mentioned that it would be even better if the other board would never see the draw offer.

\item The team Elo ratings system performed correctly for all the played games. None of the players ever had a rating below zero. The ratings were correctly scaled using K-factor. The \texttt{rating\_history} records correctly matched the records displayed in the leaderboard.

\end{itemize}

Phase 1 work could be enhanced by having more time to implement three specific features

\begin{itemize}
    \item A fully functional chat capability
    \item A takeback protocol that works well for friendly matches
    \item Tournament mode
    \item Different personalities for the two engines on a team where each engine could play a very different style and/or possess significantly different strengths.
\end{itemize}

\subsection*{8.6 Comparing Phase 1 and Phase 2}
\addcontentsline{toc}{subsection}{8.6 Comparing Phase 1 and Phase 2}

Phase 1 and Phase 2 are two completely different games.

One plays more social in Phase 1 than in Phase 2, since four players play together as teams of two players in Phase 1. As in traditional online bughouse, players communicate via an in game chat and by the timing piece transfer. The flow of captured pieces is going both ways in Phase 1. 

As mentioned in my Interim Demo, Phase 2 is independent enough that it could be a standalone variant of chess. In Phase 2 a single human player competes against a loop of chess engines. Pieces are transferred to the player’s board in one direction only, which I think creates a rather unique style of play. Because each board ends separately the games are also shorter on average, about 16 minutes as compared to about 19 minutes in Phase 1 at 5+5.

\subsection*{8.7 Limitations}
\addcontentsline{toc}{subsection}{8.7 Limitations}

\begin{enumerate}[leftmargin=*]

\item \textit{Cross-board recommendation accuracy.} While recommending strong pieces from sibling boards, the system currently is not accurate enough to account for how badly that piece weakens the remaining position on that board. A deeper multi move evaluation would help to improve this, but the cost of such evaluation increases too much for higher N values ( e.g N=10).

\item \textit{N=10 UI cramped.} The mini boards for N=10 are too small to easily read out moves and follow the resulting tactics on each board. Three other players reported this issue during testing. I added a scrollable option but the boards UI still needs to be improved.

\item \textit{High Elo engine quality deteriorates.} The high Elo Stockfish is not good enough to mimic real players with high Elo ratings. A Elo 1700+ rated engine will sometimes make obvious mistakes, but then find very strong tactical moves. This happens due to continuous change of pieces in the players reserves which continuously changes the evaluation that the engine had thought before making the previous move. Restricting the engine settings or building a proper weak engine, would be better but is out of the scope of this project.

\end{enumerate}
\subsection*{8.8 Future work}
\addcontentsline{toc}{subsection}{8.8 Future work}
\begin{itemize}

\item \textit{Better partner need logic in Phase 2.} The two functions \texttt{computePartnerNeed} and \texttt{computeOpponentPartnerDanger} (of his partner) already work fine in Phase 1, but aren’t fully connected in Phase 2. They should work fine after that as well, just need some tuning of the \texttt{LAMBDA} and \texttt{MU} values for the N boards.

\item \textit{Bughouse aware search.} As mentioned above, right now we compare the best normal move by Stockfish with the best drop as evaluated by Stockfish. A more appropriate search would be a depth 3 or 4 search including the drops in the search tree, which would likely find the strongest tactical move for a player. But this would be a major project in itself.

\item \textit{Different engine strengths per board.} Currently, all engine boards are set to the same Elo. Adding this would allow stronger players to challenge harder boards on their left and right by setting their boards to different Elo values.

\end{itemize}

\subsection*{8.9 Why Phase 2 is not rated}
\addcontentsline{toc}{subsection}{8.9 Why Phase 2 is not rated}
Infinite mode (Phase 2) does not have a rating system yet. Elo and Glicko ratings are designed with fair one vs one games in mind. The rating system of Infinite mode games is very asymmetric (one human playing against arbitary engines).
eg. An event on one board might have changes on the game as a whole, but this does not directly translate into a win or loss for the human on that board. In addition the difficulty of Infinite mode games increases with N. Thus winning an Infinite mode game at N=10 is not the same as winning a game at N=4.

There are a few ways a rating system for this could be implemented.
\begin{enumerate}
    \item Rate by engine Elo and N, so a player has a separate rating per configuration
\item Rate by average eval accuracy, engine Elo, N and the time the user lasted. I think this would give an accurate rating.
\item  Rate based on human's board outcome and not by siblings

The optimal way would require a lot of long term data that does not currently exist.
\end{enumerate}

\newpage
\section*{9. Conclusion}
\addcontentsline{toc}{section}{9. Conclusion}
I created Infinite Armada Chess in two stages. First, I created a four player bughouse game in which I wrote a server side rule enforcement layer, optionally filled out moves with Stockfish, correctly handled draws and resignations on a team basis and implemented a custom team Elo rating. Then, I extended this to an N board loop (Phase 2) where one human plays against N-1 Stockfish engines, with pieces going in one direction around the ring of boards.

The main contribution of this chess variant is the additional layer of complexity on top of normal chess namely a strict set of rules, an N board manager (currently upto N=10), a drop decision pipeline and a cross board recommendation algorithm.

This system has been tested using 150 assertion tests as well as engine strength tests and pacing tests.

This is one of the most challenging but rewarding project I have ever done. It has helped me reignite my passion for chess that I had set aside for years. The three pieces I am most proud of are the N board armada, the unidirectional cross board piece flow and the recommendation algorithm to rank sibling boards by eval cost. I had the best experience playing at N=4 to N=6. 

I was glad to come up with a version for "Infinite" myself, based on that XKCD comic. I thought it was interesting to not explore into giving my opponent infinite number of pieces, but rather redesign the topology of the game itself with connected boards and piece flow. Some chess problems are about queens, this one was about the board.

\newpage
\section*{References}
\addcontentsline{toc}{section}{References}
\begin{enumerate}[leftmargin=*]

\item Munroe, R. (2024) Infinite Armada Chess. XKCD \#3020.  \textit{\url{https://xkcd.com/3020/}} 

\item DiIorio, J. (n.d.) Bughouse Chess (Open-Source Repository). \textit{\url{https://github.com/johndiiorio/bughouse}} 

\item FIDE (2023) FIDE Laws of Chess. \textit{\url{https://handbook.fide.com/chapter/E012023}} 

\item Stockfish Chess Engine Team (n.d.) Stockfish. \textit{\url{https://stockfishchess.org/}} 

\item Lichess Organisation (n.d.) Chessground. \textit{\url{https://github.com/lichess-org/chessground}} 

\item Internet Chess Club (n.d.) Bughouse Rules. \textit{\url{https://www.chessclub.com/help/bughouse}} 

\item Glickman, M. (2012) Example of the Glicko-2 System. \textit{\url{http://www.glicko.net/glicko/glicko2.pdf}} 

\item Socket.IO Documentation (n.d.) Socket.IO Documentation. \textit{\url{https://socket.io/docs/v2/}} 

\item Supabase (n.d.) Supabase Documentation. \textit{\url{https://supabase.com/docs}} 

\item DiIorio, J. (n.d.) chess.js Fork within the Bughouse Repository, file \texttt{src/server/services/bug.js}. \textit{\url{https://github.com/johndiiorio/bughouse}} 

\item Stockfish Team (n.d.) UCI Protocol Reference. Stockfish Wiki. \textit{\url{https://github.com/official-stockfish/Stockfish/wiki/UCI}} 

\item Facebook Inc. (2017) React Version 15. \textit{\url{https://reactjs.org/blog/2017/04/07/react-v15.5.0.html}} 

\item Elo, A. E. (1978) The Rating of Chessplayers, Past and Present. New York: Arco Publishing.

\item Glickman, M. E. (1995) ‘The Glicko Rating System’, \textit{The American Statistician}.

\item Hooper, D. and Whyld, K. (1992) The Oxford Companion to Chess. 2nd edn. Oxford: Oxford University Press.

\item Pritchard, D. B. (2007) The Classified Encyclopedia of Chess Variants. Harpenden: John Beasley.

\item The Chess Variant Pages (n.d.) Bughouse Chess and Crazyhouse Entries. \textit{\url{https://www.chessvariants.com/}} 

\item Lichess Organisation (n.d.) Crazyhouse --- Variant Rules. \textit{\url{https://lichess.org/variant/crazyhouse}} 

\item Knuth, D. E. and Moore, R. W. (1975) ‘An Analysis of Alpha-Beta Pruning’, \textit{Artificial Intelligence}, 6(4), pp. 293--326.

\item Nasu, Y. (2018) NNUE: Efficiently Updatable Neural-Network Evaluation Function for Computer Shogi.

\item Silver, D., Hubert, T., Schrittwieser, J. et al. (2018) ‘A General Reinforcement Learning Algorithm that Masters Chess, Shogi, and Go through Self-Play’, \textit{Science}, 362(6419), pp. 1140--1144.

\item Pascutto, G.-C. (n.d.) Sjeng - A Computer Chess and Variants Program. \textit{\url{https://www.sjeng.org/}} 
\item Canva Pty Ltd. Canva - Online graphic design and diagramming platform.  \textit{\url{https://www.canva.com/}}  



\end{enumerate}
\newpage
\section*{Appendix A: Running the system}
\addcontentsline{toc}{section}{Appendix A: Running the system}

\textbf{GitHub Repo:} \texttt{https://github.com/devansh2605/CSCS5199-Individual-Masters-Project.git}
\textbf{Deployed for Multiplayer:} \texttt{https://infinite-chess-armada.vercel.app/}
\textbf{Diagrams:} \texttt{https://canva.link/t92g7xdswyar977}

If you want to try out the system you check the vercel deployed version or create a clone of this repo. To run the system locally it is one command \texttt{./run.sh}. This script checks whether it runs on Linux or macOS and then installs Node 20 and Stockfish locally. The script copies .env.example to .env, runs npm install, builds the webpack-bundle and starts the server. Start the app in your browser at http://localhost:1000. The.

The test suite runs with \texttt{npm test} and skips Stockfish-dependent tests cleanly if not on the path.

I have also added a \texttt{README.md} for you perusal.

\newpage
\section*{Appendix B: Testing summary}
\addcontentsline{toc}{section}{Appendix B: Testing summary}
 Section 7 described the testing methodology for this suite, as well as file by file coverage. To run the full test suite, execute \texttt{npm test}. Below is the test suite output from a complete run I performed.

\begin{verbatim}
               TESTING RESULTS
  OK   bughouse.test.js        82 pass,  0 fail,  0 skip   (114ms)
  OK   dropDecision.test.js    21 pass,  0 fail,  0 skip   (200ms)
  OK   elo.test.js             26 pass,  0 fail,  0 skip   ( 25ms)
  OK   engineStrength.test.js   2 pass,  0 fail,  0 skip   (2955ms)
  OK   infiniteMode.test.js    19 pass,  0 fail,  0 skip   (25911ms)
--------------------------------------------------
  TOTAL                   150 pass,  0 fail,  0 skip
\end{verbatim}

The three regression bugs that the tests caught during development and their fixes are described in Section 7.3. The manual testing is described in Section 7.2.

\newpage
\section*{Appendix C: Interface Screenshots}
\addcontentsline{toc}{section}{Appendix C: Interface Screenshots}

\begin{figure}[H]
    \centering
    \includegraphics[width=0.75\linewidth]{Homepage.png}
    \caption{Home Page}
\end{figure}

\begin{figure}[H]
    \centering
    \includegraphics[width=0.75\linewidth]{game-lobby.png}
    \caption{Lobby - 2vs2}
\end{figure}
\begin{figure}[H]
    \centering
    \includegraphics[width=0.75\linewidth]{phase1.png}
    \caption{Phase 1 - 2 vs 2}
\end{figure}
\begin{figure}[H]
    \centering
    \includegraphics[width=0.75\linewidth]{phase2-layout.png}
    \caption{\textit{N=4 boards in Phase 2.}}
\end{figure}

\begin{figure}[H]
    \centering
    \includegraphics[width=0.75\linewidth]{leaderboard.png}
    \caption{Leaderboard}
\end{figure}

\begin{figure}[H]
    \centering
    \includegraphics[width=0.75\linewidth]{phase2-lobby.png}
    \caption{Phase 2 - Lobby}
    
\end{figure}

\begin{figure}[H]
    \centering
    \includegraphics[width=0.75\linewidth]{50-move-draw.png}
    \caption{rare edge case - 50 move draw}
\end{figure}
\begin{figure}[H]
    \centering
    \includegraphics[width=0.75\linewidth]{login.png}
    \caption{Sign Up/Log In}
\end{figure}

\newpage
\section*{Appendix D: Other Appendices}
\addcontentsline{toc}{section}{Appendix D: Other Appendices}
\begin{enumerate}
    \item Ethics Application
    \item Description, Objectives, Ethics and Resources
    \item Plan \& Context Survey
    \item Interim Demo
\end{enumerate}
\begin{figure}[H]
    \centering
    \includegraphics[width=0.75\linewidth]{ethics-application.png}
\end{figure}
\end{document}

